\section{Outros Modelos Contínuos}

\subsection{Uniforme Contínua}

Analogamente à distribuição uniforme discreta, na distribuição uniforme contínua, todos
os valores no suporte possuem a mesma densidade de probabilidade.
Os parâmetros dessa distribuição definem o suporte dessa variável, e a densidade
de probabilidade sobre cada valor possível é o inverso do comprimento do intervalo.

\begin{def:uniforme continua}
Uma variável $X$ tem distribição uniforme contínua com parâmetros $a$ e $b$ caso a sua
função de densidade seja dada por
\[
    f(x) = \begin{cases}
    \frac{1}{b-a},&\text{se } a \le x \le b \\
    0,&\text{caso contrário}
    \end{cases}
\]
\end{def:uniforme continua}

\begin{prop:uniforme continua}
\label{prop:uniforme continua}
Se $X \sim \mathrm{Unif}(a, b)$, então
\begin{enumerate}
\item sua função de distribuição é dado por
\[
    F(x) = \begin{cases}
        0,&\text{se } x < a \\
        \frac{x - a}{b-a},&\text{se } a \le x \le b \\
        1,&\text{se } x > b \\
    \end{cases}
\]
\item $\ev{X} = \frac{b + a}{2}$
\item $\var{X} = \frac{(b-a)^2}{12}$
\end{enumerate}
\end{prop:uniforme continua}

\begin{obs:uniforme continua 1}
A função de distribuição da variável uniforme contínua cresce linearmente a partir de
$(a, 0)$ até $(b,1)$.
Portanto, a probabilidade de um determinado intervalo $(\alpha, \beta) \subset (a, b)$
é proporcional ao tamanho do próprio intervalo:
\[
    F(\alpha \leq X \leq \beta) = F(\beta) - F(\alpha) = \frac{\beta - \alpha}{b - a} \propto (\beta - \alpha)
\]
\end{obs:uniforme continua 1}

A próxima proposição mostra que uma determinada transformação aplicada numa variável
uniforme sempre resulta na mesma variável uniforme, independente dos parâmetros da
distribuição original.

\begin{prop:transformacao uniforme continua}
Se $X \sim \mathcal{U}(a, b)$ e
\[
    U = \frac{X - \frac{a+b}{2}}{b-a}
\]
então $U \sim \mathcal{U}(-1/2, 1/2)$.
\end{prop:transformacao uniforme continua}
\begin{proof}
Se $u = g(x)$ é a transformação dada por
\[
    \frac{x - \frac{a+b}{2}}{b-a},
\]
então ela é uma transformação monótona crescente.
Pelo Teorema \ref{thm:fda transformacao continua},
\[
   F_U(u) = F_X(g^{-1}(u))
\]
onde $F_U(u)$ e $F_X(x)$ são as funções de distribuição de $U$ e $X$ respectivamente.
Se $a \le x \le b$, então
\[
   F_U(u) = \frac{1}{b-a} [u(b-a) + (b+a)/2] - \frac{a}{b-a},
\]
e, após algumas manipulações algébricas, chegamos a
\[
    F_U(u) = \frac{(b-a)u + \frac{b-a}{2}}{b-a} = u + \frac{1}{2}.
\]
Para determinar os parâmetros $\alpha$ e $\beta$ da distribuição de $U$, percebamos que
\[
    \begin{cases}
        \beta-\alpha = 1 \\
        -\alpha = 1/2
    \end{cases}
\]
cuja solução é $\alpha = -1/2$ e $\beta = 1/2$, logo $U \sim \mathcal{U}(-1/2, 1/2)$.
\end{proof}

Essa Proposição é útil para que programas de computador, por exemplo, possam tabelar
valores para uma única distribuição uniforme e, com isso, sejam capazes de determinar
valores de probabilidade para qualquer outra distribuição uniforme de forma eficiente.

\subsection{Exponencial}

A variável exponencial é rotineiramente aplicada para modelar tempos de espera entre
eventos os quais ocorrem a uma taxa fixa, como chamadas telefônicas a um
\textit{call center} ou de requisições a um servidor de dados.
Sua função de densidade inclui um parâmetro $\lambda$, que denota a taxa de ocorrência
de eventos, de sorte que o seu inverso denota o tempo médio entre dois eventos.

\begin{def:exponencial}
A variável $X$ tem distribuição exponencial com parâmetro $\lambda > 0$ se sua função de
densidade for dada por
\[
    f(x) = \begin{cases}
        \lambda e^{-\lambda x},&\text{se } x \ge 0 \\
        0,&\text{caso contrário.}
    \end{cases}
\]
\end{def:exponencial}

\begin{prop:exponencial}
Se $X \sim \mathrm{Exp}(\lambda)$, então
\begin{enumerate}
\item Sua função de distribuição é dada por
\[
    F(x) = \begin{cases}
        1 - e^{\lambda x},&\text{se } x \ge 0 \\
        0,&\text{caso contrário.}
    \end{cases}
\]
\item $\ev{X} = \frac{1}{\lambda}$
\item $\var{X} = \frac{1}{\lambda^2}$
\end{enumerate}
\end{prop:exponencial}

\begin{obs:exponencial 1}
A variável eponencial é a contraparte contínua da variável geométrica quanto a propriedade
de sem memória.
Basta observar que
\[
    P(X \ge s \mid X \ge t) =
    \frac{P(X \ge s)}{P(X \ge t)} =
    \frac{1 - P(X \le s)}{1 - P(X \le t)} =
    \frac{e^{-\lambda s}}{e^{-\lambda t}} =
    P(X \ge s - t).
\]
É verdade também que a variável exponencial é a única dentre as contínuas com tal propriedade
\end{obs:exponencial 1}

\subsection{Distribuição Gama}

A variável com distribuição $\Gamma$ é uma generalização da distribuição exponencial.
A integral
\[
   \int_0^\infty t^{\alpha - 1} e^{-t}\, dt
\]
denota a função $\Gamma(\alpha)$, a qual é finita quando $\alpha > 0$, podendo ser
expressa por uma fórmula fechada.
Essa função satisfaz algumas relações importantes; em particular
\[
    \Gamma(\alpha+1) = \alpha \Gamma(\alpha),\, \alpha > 0.
\]
Uma vez que $\Gamma(1) = 1$, então, pela relação anterior, vale que $\Gamma(n) = n!$ para
todo $n$ inteiro positivo.

Como o integrando na definição da função $\Gamma$ é positivo, então segue de imadiato
que
\begin{equation}
\label{eq:pdf gamma simples}
    f(t) = \frac{t^{\alpha-1}e^{-t}}{\Gamma{\alpha}}
\end{equation}
é uma função de densidade.
No entanto, é possível generalizar a distribuição com uma parâmetro adicional $\beta$.
Se $T$ tem densidade igual àquela dada por \ref{eq:pdf gamma simples} e $X = \beta T$,
então defini-se a distribuição $\Gamma$ da maneira a seguir:

\begin{def:gama}
A variável $X$ tem distribuição $\Gamma$ com parâmetros positivos $\alpha$ e $\beta$ se
sua função de densidade for dada por
\[
    f(x) = \frac{1}{\Gamma(\alpha) \beta^\alpha} x^{\alpha - 1} e^{-x / \beta}
\]
onde a função $\Gamma(\alpha)$ é definida como
\[
    \Gamma(\alpha) = \int_0^\infty t^{\alpha - 1} e^{-t}\, dt.
\]
\end{def:gama}

O parâmetro $\alpha$ é chamado de \textbf{forma} (\textit{shape}), pois caracteriza o pico
da distribuição, enquanto que o parâmetro $\beta$ é denominado \textbf{escala}
(\textit{scale}), já que influencia a dispersão dela.
Percebamos também que a distribuição exponencial de taxa $\lambda$ é um caso particular da
distribuição $\Gamma$, onde $\alpha = 1$ e $\beta = 1/\lambda$.
Com efeito,
\[
    f(x \mid \alpha = 1, \beta = 1/\lambda) = \frac 1 \beta e^{-x/\beta} = \lambda e^{-\lambda x}
\]

\begin{prop:gama}
Se $X \sim \Gamma(\alpha, \beta)$, então
\[
    \ev{X} = \alpha \beta \quad \text{e} \quad \var{X} = \alpha \beta^2
\]
\end{prop:gama}

\begin{prop:mgf gama}
Se $t < 1/\beta$, então a função geradora de momentos da distribuição $\Gamma$ está
definida e é determinada por
\[
    M_X(t) = \left(\frac{1}{1 - \beta t}\right)^\alpha
\]
\end{prop:mgf gama}

Assim como a exponencial é a contraparte contínua da variável geométrica, e
a variável $\Gamma$ é a contraparte contínua da variável binomial negativa.
Com isso, a distribuição $\Gamma$ é compreendida como a soma de variáveis exponenciais, e
pode ser usada para modelar contagem de eventos até um número desejado $\alpha$ que
ocorrem com tempo médio de intervalos iguais a $\beta$.

\begin{prop:relacao gama e poison}
Seja $X \sim \Gamma(\alpha, \beta)$.
Para todo $x \in \real$, se $Y \sim \mathrm{Poisson}(x/\beta)$, então
\[
    P(X \le x) = P(Y \ge \alpha)
\]
\end{prop:relacao gama e poison}

\subsection{Distribuição $\chi^2$}

Outra especialização da distribuição $\Gamma$ é a $\chi^2$ com $n$ graus de liberdade,
onde $\alpha = n/2$ e $\beta=2$.
Ela é uma das principais distribuiçõe utilizadas em inferência estatística, sendo usada
para teste de signficância estatística de parâmetros de outras distribuições e para
amostragem da distribuição normal.

\begin{def:chi quadrado}
A variável $X$ tem distribuição $\chi^2$ com $n$ graus de liberdade se sua função de
densidade for dada por
\[
    f(x) = \begin{cases}
        \frac{1}{2^{n/2} \Gamma(n/2)} x^{n/2}e^{-x/2},&\text{se } x \ge 0 \\
        0,&\text{caso contrário.}
    \end{cases}
\]
\end{def:chi quadrado}

\begin{prop:chi quadrado}
Se $X \sim \chi^2(n)$, então
\[
    \ev{X} = n \quad \text{e} \quad \var{X} = 2n
\]
\end{prop:chi quadrado}

\begin{prop:chi quadrado e normal}
Seja $(Z_i)_{i=1}^n$ uma sequência de variáveis aleatórias com distribuição
$\mathrm{Normal}(0, 1)$.
A variável
\[
    X = \sum_{i=1}^n Z_i^2
\]
tem distribuição $\chi^2$ com $n$ graus de liberdade.
\end{prop:chi quadrado e normal}

\subsection{Ditribuição de Cauchy}

A distribuição de Cauchy, assim como a normal, é simétrica e possui um forato de sino,
porém com caldas mais pesadas.
Sua função é definida a seguir:

\begin{def:cauchy}
Uma variável aleatória $X$ tem distribuição de Cauchy com parâmetros $\theta \in \real$ e
$\gamma > 0$ se sua função de densidade for
\[
    f(x) = \frac{1}{\pi\gamma} \left[1 + \left(\frac{x - \theta}{\gamma}\right)^2\right]^{-1}
\]
\end{def:cauchy}

Uma propriedade dessa distribuição é que seus momentos não existem, uma vez que todos
eles são infinitos.
Por consequência, a função geradora de momentos dessa variável não é definida.

\begin{prop:cauchy nao tem momentos}
Para todo $\theta \in \real$, se $X \sim \mathrm{Cauchy}(\theta)$, então
\[
    \ev{X^n} = \infty
\]
para todo $n$ inteiro positivo.
\end{prop:cauchy nao tem momentos}

Apesar da inexistência de momentos, o formato da distribuição com pico $\theta$ sugere
um fato interessante.
Como visto, a distribuição de Cauchy não possui média, e, na verdade, $\theta$ denota
a mediana da distribuição, o que é enunciado na próxima Proposição.

\begin{prop:mediana de cauchy}
Se $X \sim \mathrm{Cauchy}(\theta)$, então
\[
    P(X \le \theta) = \frac 1 2
\]
\end{prop:mediana de cauchy}

A distrbuição de Cauchy é muito utilizada para modelar razões entre duas outras variáveis.
Em verdade, é possível mostrar que a razão entre duas variáveis com distribuição normal
padrão tem distribuição de Cauchy.

\subsection{Distribuição de Weilbull}

A variável de Weilbull é um modelo de distribuição ainda mais geral que a distribuição
$\Gamma$ usada para modelar tempos de vida. Além dos parâmetros de forma e escala,
adiciona-se mais um parâmetro de localização, inidicando o tempo mínimo antes de uma
falha poder ocorrer.

\begin{def:weilbull}
A variável $X$ tem distribuição de Weibull com parâmetros positivos $\alpha$, $\beta$ e
$\gamma$ se sua função de densidade for dada por
\[
    f(x) = \begin{cases}
        \frac \beta \alpha \left(\frac{x - \gamma}{\alpha}\right)^{\beta - 1} e^{-\left(\frac{x - \gamma}{\alpha}\right)^{\beta}},&\text{se } x \ge \gamma \\
        0,&\text{caso contrário.}
    \end{cases}
\]
\end{def:weilbull}

Uma interpretação pode ser atribuía a cada parâmetro dessa distribuição:
\begin{itemize}
\item Parâmetro de forma $\alpha$ -- Se $\alpha < 1$, então a taxa de falha diminui com o
tempo, a exemplo da taxas de mortalidade infantil com relação a idade.
Caso $\alpha = 1$, então as falhas ocorrem de maneira aleatória em todo período de tempo
observado, similar à distribuição exponencial.
Já se $\alpha > 1$, então a taxa de falha aumenta com o tempo, como peças que se desgatam
com o tempo.

\item Parâmetro de escala $\beta$ -- modela a dispersão da distribuição no eixo do tempo,
e, a medida que cresce, mais probabilidade é atribuída à valores maiores.
Seu argumento representa o quantil no qual cerca de $63.2\%$ da população terá apresentado
alguma falha.

\item parâmetro de localização $\gamma$ -- indica o tempo mínimo antes que seja possível
de ocorrer uma falha.
Corresponde a origem da curva de densidade no eixo do tempo.
\end{itemize}

\subsection{Distribuição Beta}

A distribuição Beta possui a propriedade de concentrar toda probabilidade num único
intervalo, sendo esse, nesse caso, o intervalo $(0, 1)$.
Com isso, ela é comumente aplicada para modelar proporções e taxas que se encontram
no intervalo citado.

\begin{def:beta}
A variável $X$ tem distribuição beta com parâmetros positivos $a$ e $b$ se sua densidade
for dada por
\[
    f(x) = \begin{cases}
        \frac{1}{B(a, b)} x^{a-1} (1-x)^{b-1},&\text{se } 0 < x < 1 \\
        0,&\text{caso contrário.}
    \end{cases}
\]
onde
\[
    B(a, b) = \int_0^1 x^{a - 1} (1-x)^{b-1}\,dx = \frac{\Gamma(\alpha)\Gamma(\beta)}{\Gamma(\alpha+\beta)}
\]
\end{def:beta}

Essa distribuição também possui a característica de assumir formas totalmente distintas
a depender dos valores dos parâmetros.
Em síntese, temos que
\begin{itemize}
\item $\alpha > 1$ e $\beta = 1$ -- distribuição estritamente crescente.
\item $\alpha = 1$ e $\beta > 1$ -- distribuição estritamente decrescente.
\item $\alpha < 1$ e $\beta < 1$ -- formato de ``U''.
\item $\alpha > 1$ e $\beta > 1$ -- formato unimodal.
\item $\alpha = \beta$ -- simétrica entorno de $1/2$ e com variância $[4(2\alpha+1)]^{-1}$.
\end{itemize}

No último caso, a distribuição torna-se cada vez mais concentrada a medida que $\alpha$
cresce.
Por fim, quando $\alpha = \beta = 1$, a distribuição degenera-se na uniforme no intervalo
$(0, 1)$.

\subsection{Distribuição log-normal}

Quando lidando com dados assimétricos à direita (cujos distribuição concentra-se à
esquerda da mediana), é muito comum aplicar a transformação $\log(x)$ para corrigir
essa assimetria, pois os valores na calda à direita são ``puxados'' para estarem
mais próximos dos menores.
Quando essa transformação leva a uma distribuição normal, então dizemos que a variável
original tem distribuição log-normal.

\begin{def:lognormal}
A variável $X$ tem distribuição log-normal com parâmetros $\mu$ e $\sigma^2$ se
sua função de densidade for dada por
\[
    f(x | \mu, \sigma^2) = \frac{1}{\sqrt{2\pi\sigma^2}} \frac 1 x e^{-(\log x - \mu)^2/(2\sigma^2)}
\]
\end{def:lognormal}

\begin{prop:lognormal}
Se $\log X \sim \mathrm{Normal}(\mu, \sigma^2)$, então
\[
    \ev{X} = e^{\mu + (\sigma^2/2)} \quad \text{e} \quad \var{X} = e^{2(\mu+\sigma^2)} - e^{2\mu + \sigma^2}
\]
\end{prop:lognormal}

\subsection{Distribuição de Laplace}

A distribuição de Laplace é obtida ao espelhar a distribuição exponencial em torno da
sua média.
Visualmente, ela possui caldas mais pesada e longas do que a normal.
Além disso, a curva da distribuição no ponto $x = \mu$ é um ``bico'', ou seja, a
densidade não é diferenciável naquele ponto.

\begin{def:laplace}
A variável $X$ tem distribuição de Laplace se sua função de densidade é dada por
\[
    f(x \mid \mu, \sigma) = \frac{1}{2\sigma} e^{-|x - \mu|/\sigma}
\]
\end{def:laplace}

Uma das aplicações da distribuição de Laplace é na seleção de atributos para uma
regressão linear, atribuído pesos iguais a zero às características consideras não
significativas para explicar outra variável.
